John R. Fanchi's Parameterized Relativistic Dynamics (PRD)
Fanchi's work is much closer to the Stueckelberg–Horwitz–Piron (SHP) programme that we have been discussing.

My assessment is that Fanchi's PRD belongs to a small, sustained research programme in manifestly covariant relativistic mechanics. It has a legitimate publication history, including papers in Physical Review A, Foundations of Physics and other journals, but it has not been adopted as a mainstream alternative to conventional relativistic quantum mechanics or quantum field theory.

There is an important distinction, however: Fanchi's work is not simply a reformulation of SHP. He has tried to give parametrized relativistic quantum mechanics its own probabilistic and physical foundations.
1. What exactly is Fanchi proposing?
The central idea is an invariant evolution parameter \(s\), independent of the spacetime coordinates \(x^\mu\).

The quantum state evolves according to an equation of the form
\[
i\hbar\frac{\partial\psi}{\partial s}
=K\psi.
\]

Here:
- \(x^\mu=(ct,\mathbf x)\) are spacetime coordinates.
- \(s\) is an invariant evolution parameter.
- \(K\) is the generator of translations in \(s\).
- The particle is not required to remain on a fixed mass shell at every instant of evolution.
This last point is essential. The theory permits off-shell dynamics.

Fanchi's early papers make a particularly interesting distinction between three notions of time:
1. Galilean time: the evolution parameter of nonrelativistic mechanics.
2. Minkowski time: the temporal coordinate of spacetime.
3. Historical time: the invariant ordering parameter of parametrized relativistic mechanics.
His 1986 paper explicitly develops an operational procedure for measuring the evolution parameter using a physical clock. 
Phys. Rev. A+1

2. The scientific publication history
Fanchi has developed this programme over several decades.

Some particularly relevant publications are:
Year	Publication	Subject
1978	Foundations of Physics	Quantum mechanics of relativistic spinless particles
1986	Physical Review A	Parametrizing relativistic quantum mechanics
1988	Physical Review A	Gibbs ensemble and cosmological implications
1993	Foundations of Physics	Review of invariant-time formulations
1994	Foundations of Physics	Validity of parametrized relativistic wave equations
2000	Foundations of Physics	Quantum potential in relativistic dynamics
2008	Hadronic Journal	Tutorial on Parameterized Relativistic Dynamics
2011	Foundations of Physics	Manifestly covariant quantum theory
2019	Symmetry	Neutrino flavor transitions
2025	Mathematics	Probabilistic formulation in curved spacetime

Sources: 
PubMed+7


This is a substantial publication history, although the number of publications should not be confused with the degree of independent acceptance.

An especially important paper is his 1994 Foundations of Physics article, Evaluating the Validity of Parametrized Relativistic Wave Equations.

Fanchi examines possible evolution generators and argues that the generator proportional to the four-momentum squared,
\[
K\propto P^\mu P_\mu,
\]

is the one justified by his probabilistic construction. He nevertheless explicitly acknowledges that experimental testing is necessary. 
PhilPapers


That qualification is important: even within the programme, the selection of the evolution generator is not simply a mathematical inevitability.
3. How does the scientific community regard it?
I would characterize the situation as follows.

A. It is a recognized alternative formulation.

Fanchi's work has appeared in established physics journals. The invariant-parameter approach is also discussed by other researchers, particularly Matej Pavšič and Lawrence Horwitz.

B. It has not become an independently established physical theory.

The central ideas are known, but PRD has not acquired a substantial independent research community comparable to quantum field theory or mainstream relativistic quantum mechanics.

C. Its most important claims remain open to experimental assessment.

Fanchi has proposed applications to:
- Neutrino oscillations.
- Particle decay.
- Relativistic bound states.
- Nonlocality.
- Tachyonic motion.
For example, his 2021 paper treats exponential particle decay and neutrino oscillations within parametrized relativistic quantum mechanics. 
repository.tcu.edu


These are interesting applications, but proposing a different mathematical description of an observed phenomenon is not equivalent to demonstrating a new physical effect.
4. Comparison with SHP
This is perhaps the most important issue for our purposes.

Fanchi's PRD and SHP share the same broad conceptual architecture:
Feature	Fanchi PRD	SHP
Evolution parameter	Invariant \(s\)	Invariant \(\tau\)
Spacetime coordinates	\(x^\mu\)	\(x^\mu\)
Off-shell dynamics	Yes	Yes
Manifest covariance	Yes	Yes
Quantum evolution	Schrödinger-like equation in \(s\)	Stueckelberg–Horwitz–Piron equation
Electromagnetic interactions	Gauge-field formulation	Five-dimensional gauge structure
Many-particle theory	Yes	Yes
Curved spacetime	Developed in later work	Developed by Horwitz and collaborators

The distinction is not absolute. Both belong to the broader parametrized relativistic quantum theory tradition.

However, I would emphasize three differences.

First: physical foundations.

Fanchi places considerable emphasis on the probabilistic foundations of the theory, including the meaning and measurement of the invariant evolution parameter.

Second: the mass operator.

Fanchi investigates the relationship between the evolution generator and the mass spectrum. This is particularly relevant to his treatment of neutrino oscillations as transitions between mass states.

Third: interpretation.

Fanchi explicitly presents PRD as an alternative to the conventional QFT paradigm. SHP, in contrast, has developed into a broader research programme involving classical mechanics, electrodynamics, statistical mechanics and gauge theory.

The comparison is not entirely symmetrical: SHP has a larger independent research literature.
5. What are the most serious scientific questions?
I would not begin by asking whether PRD is mathematically consistent. I would ask whether its additional structure produces sufficiently compelling physics.

Here are four questions I consider decisive.

1. The meaning of \(s\). Is the invariant evolution parameter a new physical time, an ordering parameter, or an auxiliary parameter? Can its measurement be connected unambiguously to actual clock readings?

2. The mass-shell problem. What determines the distribution of invariant masses? How does the theory recover the observed mass spectrum of stable particles?

3. The relationship with QFT. Can PRD reproduce particle creation and annihilation, scattering amplitudes, and the quantum statistics of identical particles without effectively reintroducing QFT?

4. Distinctive predictions. Are there experimentally accessible situations in which PRD predicts something different from conventional relativistic quantum mechanics and QFT?

The fourth question is particularly important. Neutrino oscillations are a promising subject for comparison, but the standard three-flavor framework already describes a large body of experimental data. A competing theory must be tested against the same data, with its parameters and uncertainties properly accounted for.
6. What I would recommend reading
For a serious examination, I would choose these four papers:
1. Fanchi (1993)

Review of Invariant Time Formulations of Relativistic Quantum Theories.
Foundations of Physics, 23, 487–548 (1993)Start here
Particularly valuable for understanding the wider family of invariant-time theories.


2. Fanchi (1994)

Evaluating the Validity of Parametrized Relativistic Wave Equations.
Foundations of Physics, 24, 543–562 (1994)Important for understanding the choice of evolution generator.


3. Fanchi (2011)

Manifestly Covariant Quantum Theory with Invariant Evolution Parameter in Relativistic Dynamics.
Foundations of Physics, 41, 4–32 (2011)A useful systematic presentation of his formulation.


4. Fanchi (2025)

Probabilistic Basis of Parametrized Relativistic Quantum Theory in Curved Spacetime.
Mathematics, 13, 1657 (2025)Particularly relevant to extending parametrized dynamics to general relativity.



7. My conclusion for our research
I would take Fanchi seriously, but distinguish his achievements from the claims made for PRD as an alternative physical paradigm.

In particular, I think his work deserves examination in connection with three issues we have already discussed:
- The status of global evolution time in relativistic mechanics.
- The distinction between on-shell and off-shell dynamics.
- The possibility of constructing a genuinely relativistic stochastic mechanics.
There is one further point I find especially interesting.

Fanchi's invariant evolution parameter and Nelson's stochastic time need not have the same physical meaning.

Nelson's time is used to define forward and backward diffusion processes. Fanchi's \(s\) orders the evolution of relativistic quantum states.

It would be interesting to investigate whether the latter can provide the evolution parameter required by a relativistic Nelson theory, without imposing the usual mass-shell condition at every intermediate stage.

That would connect three programmes:
\[
\boxed{
\text{Stueckelberg–Horwitz–Piron}
\longrightarrow
\text{Fanchi PRD}
\longrightarrow
\text{Relativistic Nelson mechanics}
}
\]

I would regard this as a worthwhile research direction, but not assume that the connection works until its probability conservation, stochastic dynamics, and nonrelativistic limit have been established.

One suggestion: before attempting the Nelson extension, I would like to examine Fanchi's 1994 paper in detail, especially his probabilistic argument for selecting \(K\propto P^\mu P_\mu\). That is arguably the most consequential foundational step in his formulation.
Fanchi (1994): Evaluating the Validity of Parametrized Relativistic Wave Equations
Let us examine the argument carefully, because I think this paper contains one of the most consequential—and debatable—steps in Fanchi's entire programme.

The question is not merely why one should introduce an invariant evolution parameter \(s\). It is why the generator of evolution in \(s\) should be quadratic in the four-momentum.

Fanchi argues that a probabilistic foundation selects this generator over competing possibilities.

My conclusion, which I will develop below, is that his argument gives a physically motivated reason for preferring the quadratic generator, but does not establish its uniqueness from probability theory alone.

There are three logically distinct steps:
1. Introduce a probability distribution over spacetime, conditional on \(s\).
2. Derive a probability conservation equation.
3. Construct a wave equation compatible with that conservation equation and identify its evolution generator.
The distinction between steps 2 and 3 is especially important.

The original reference

J. R. Fanchi, Evaluating the Validity of Parametrized Relativistic Wave Equations, Foundations of Physics 24, 543–562 (1994).

1. The starting point: probability in spacetime
In conventional nonrelativistic quantum mechanics, we have
\[
\rho(\mathbf x,t)=|\psi(\mathbf x,t)|^2
\]

and the normalization condition
\[
\int_{\mathbb R^3}\rho(\mathbf x,t)\,d^3x=1.
\]

Notice that the normalization is performed over space at a fixed time.

In a manifestly covariant relativistic theory, Fanchi proposes a different construction:
\[
\boxed{\rho(x|s)=|\psi(x,s)|^2}
\]

where
\[
x^\mu=(ct,\mathbf x)
\]

and \(s\) is an invariant evolution parameter.

The normalization becomes
\[
\boxed{
\int_{\mathbb R^4}\rho(x|s)\,d^4x=1.
}
\tag{1}
\]

This is the central conceptual change.

At a given \(s\), the particle has a probability distribution over spacetime, not merely over three-dimensional space.

For example, consider
\[
\rho(x|s)=|\psi(x,s)|^2.
\]

The probability of detecting the particle in a spacetime region \(D\), conditional on the value of \(s\), is
\[
P(D|s)=\int_D \rho(x|s)\,d^4x.
\]

This permits the probability of finding a particle at a particular coordinate time to change with \(s\).

In particular, the marginal distribution in the temporal coordinate is
\[
\rho_t(t|s)=
\int_{\mathbb R^3}\rho(ct,\mathbf x|s)\,d^3x.
\]

It is therefore possible for the particle to have a finite lifetime: its probability of being detected at a given coordinate time can become negligible.

This is one of the motivations for Fanchi's approach to particle stability and decay. 
Ruggero Maria Santilli Foundation+1


There is an important conceptual qualification: this is a probability distribution on four-dimensional spacetime, conditioned on \(s\). It is not the probability that the particle occupies all spacetime simultaneously.
2. Probability conservation
Since the total probability must remain equal to one as \(s\) evolves,
\[
\frac{d}{ds}\int_{\mathbb R^4}\rho(x|s)\,d^4x=0.
\]

Assuming sufficient decay at spacetime infinity, this implies a local conservation equation.

Fanchi introduces a four-dimensional probability flux
\[
J^\mu(x,s)=\rho(x|s)V^\mu(x,s),
\]

where \(V^\mu\) is interpreted as a four-velocity field.

The continuity equation is
\[
\boxed{
\frac{\partial\rho}{\partial s}
+
\partial_\mu J^\mu=0.
}
\tag{2}
\]

Equivalently,
\[
\boxed{
\frac{\partial\rho}{\partial s}
+
\partial_\mu(\rho V^\mu)=0.
}
\tag{3}
\]

This is the relativistic counterpart of the ordinary continuity equation
\[
\partial_t\rho+\boldsymbol\nabla\cdot\mathbf j=0.
\]

There is a subtle but important difference.

In ordinary quantum mechanics, the continuity equation describes conservation of probability in three-space as time evolves.

Here it describes conservation of probability in four-dimensional spacetime as the invariant parameter \(s\) evolves.

The new evolution equation is therefore a five-dimensional continuity equation, with \(s\) playing the role of an evolution variable rather than a fifth spacetime coordinate.
3. Introducing the wavefunction
Fanchi adopts the Born representation
\[
\rho=\psi^*\psi.
\tag{4}
\]

He also introduces a phase representation
\[
\psi(x,s)=\sqrt{\rho(x,s)}
\exp\left(\frac{i}{\hbar}S(x,s)\right).
\tag{5}
\]

This is useful for deriving the relationship between probability flow and momentum.

For electromagnetic interactions, define the kinetic four-momentum operator
\[
\hat\pi_\mu=
-i\hbar\partial_\mu-\frac ec A_\mu
\]

(using the convention in which the particle charge is \(e\)).

The corresponding probability current is taken to be
\[
J^\mu=
\frac{1}{2m}
\left[
\psi^*\hat\pi^\mu\psi+
(\hat\pi^\mu\psi)^*\psi
\right].
\tag{6}
\]

Consequently,
\[
V^\mu=\frac{J^\mu}{\rho}.
\]

Substituting the polar representation gives
\[
\boxed{
V^\mu=
\frac{1}{m}
\left(
\partial^\mu S-\frac ec A^\mu
\right).
}
\tag{7}
\]

This resembles the familiar Hamilton–Jacobi relation between momentum and the gradient of the action.

It is also closely related to the velocity field in Nelson's mechanics and in the Madelung representation.

However, we should not confuse this with a derivation of Nelson mechanics: there is no diffusion process or pair of forward and backward derivatives here.
4. The crucial derivation of the wave equation
Now substitute (4) and (6) into the continuity equation (2).

One obtains
\[
\frac{\partial}{\partial s}|\psi|^2
+
\partial_\mu
\left[
\frac{1}{2m}
\left(
\psi^*\hat\pi^\mu\psi+
(\hat\pi^\mu\psi)^*\psi
\right)
\right]=0.
\tag{8}
\]

Using the fact that the electromagnetic covariant derivative is Hermitian with respect to the appropriate spacetime inner product, this can be written as
\[
\psi^*
\left[
\frac{\partial\psi}{\partial s}
-\frac{i\hbar}{2m}
D_\mu D^\mu\psi
\right]
+
\text{complex conjugate}=0,
\]

where
\[
D_\mu=\partial_\mu+
\frac{ie}{\hbar c}A_\mu.
\]

A sufficient way to satisfy this equation is to impose
\[
\boxed{
i\hbar\frac{\partial\psi}{\partial s}
=
\frac{1}{2m}\hat\pi_\mu\hat\pi^\mu\psi.
}
\tag{9}
\]

With an additional scalar interaction potential \(V\), Fanchi writes the generator in the form
\[
\boxed{
K=\frac{1}{2m}\hat\pi_\mu\hat\pi^\mu+V.
}
\tag{10}
\]

Thus,
\[
\boxed{
i\hbar\partial_s\psi=K\psi.
}
\tag{11}
\]

This is the Stueckelberg-type parametrized wave equation.

The free-particle equation, in the absence of electromagnetic and scalar interactions, is
\[
i\hbar\frac{\partial\psi}{\partial s}
=
-\frac{\hbar^2}{2m}
\partial_\mu\partial^\mu\psi.
\tag{12}
\]

Fanchi's 2008 tutorial explicitly presents this probabilistic construction and the resulting mass operator. 
Ruggero Maria Santilli Foundation+1

5. Why does Fanchi prefer the quadratic generator?
This is the actual subject of the 1994 paper.

Fanchi compares three proposed generators.
A. Dirac-like generator
\[
K_1=\alpha_1\gamma^\mu P_\mu.
\tag{13}
\]

This is a first-order differential operator.

It has the advantage of resembling the Dirac equation and can accommodate spinorial wavefunctions.

However, Fanchi argues that the associated probability normalization and current are problematic in the parametrized four-dimensional formulation.
B. Quadratic generator
\[
\boxed{
K_2=\alpha_2 P_\mu P^\mu.
}
\tag{14}
\]

This is the Fock–Stueckelberg choice.

It is a second-order differential operator and has a direct classical counterpart.

With electromagnetic interaction, it becomes
\[
K_2=
\frac{1}{2m}
\left(P_\mu-\frac ec A_\mu\right)
\left(P^\mu-\frac ec A^\mu\right).
\tag{15}
\]

This generator admits the positive-definite spacetime probability density and the associated conventional probability current.

It also connects naturally to the classical relativistic Hamiltonian.
C. Square-root generator
\[
K_3=\alpha_3\sqrt{P_\mu P^\mu}.
\tag{16}
\]

This is another possible choice, related to a relativistic mass operator.

It introduces difficulties associated with defining the square root, especially for spacelike momenta and an unrestricted off-shell spectrum.

Fanchi's assessment is that \(K_2\) is the most satisfactory choice because it admits a straightforward probabilistic formulation and classical correspondence.

The 1994 abstract states his conclusion quite explicitly: \(K_2\) is the only one of the proposed generators that he can derive from first principles, notably a probabilistic basis. He nevertheless says that experimental tests are needed. 
Springer Nature Link+1


6. The connection with classical relativistic mechanics
I think this is an especially important part of the argument.

Consider the classical Hamiltonian
\[
K(x,p)=\frac{1}{2m}p_\mu p^\mu.
\tag{17}
\]

Hamilton's equations with respect to \(s\) are
\[
\frac{dx^\mu}{ds}
=
\frac{\partial K}{\partial p_\mu}
=
\frac{p^\mu}{m},
\tag{18}
\]

and
\[
\frac{dp_\mu}{ds}
=
-\frac{\partial K}{\partial x^\mu}.
\tag{19}
\]

For a free particle,
\[
\frac{dp_\mu}{ds}=0.
\]

Therefore,
\[
p^\mu=m\frac{dx^\mu}{ds}.
\tag{20}
\]

This is precisely the familiar relativistic momentum-velocity relation when \(s\) is identified with proper time on a timelike trajectory.

But there is a crucial difference.

In ordinary relativistic mechanics, the mass-shell condition is imposed:
\[
p_\mu p^\mu=m^2c^2.
\tag{21}
\]

In parametrized dynamics, this condition need not be imposed as an identity on the dynamical phase space.

Instead, the value of \(p^2\) determines the value of the evolution generator.

For the free particle,
\[
K=\frac{p^2}{2m}.
\]

On the conventional mass shell,
\[
K=\frac{mc^2}{2}.
\]

But off shell, \(K\) can take other values.

This is the fundamental source of the additional dynamical possibilities.
A useful connection with Hamilton–Jacobi theory
The corresponding classical Hamilton–Jacobi equation is
\[
\frac{\partial S}{\partial s}
+
\frac{1}{2m}
\left(
\partial_\mu S-\frac ec A_\mu
\right)
\left(
\partial^\mu S-\frac ec A^\mu
\right)
+V=0.
\tag{22}
\]

Notice that this is an ordinary Hamilton–Jacobi equation, but with four spacetime coordinates treated as configuration variables and \(s\) as evolution parameter.

This is one reason I find Fanchi's approach interesting in connection with our earlier discussions of Hamilton–Jacobi theory and parametrized mechanics.

It avoids treating coordinate time as the unique evolution variable.
7. The most important question: does probability conservation really determine \(K_2\)?
Here I would be rather more critical than Fanchi.

His argument is physically appealing, but there is a logical gap between deriving a continuity equation and uniquely determining the generator.

Let us examine it.

We have
\[
\partial_s\rho+\partial_\mu J^\mu=0.
\tag{23}
\]

We also have
\[
\rho=\psi^*\psi.
\]

Suppose we choose the current
\[
J^\mu=
\frac{1}{2m}
\left[
\psi^*(-i\hbar D^\mu\psi)
+
(i\hbar D^\mu\psi)^*\psi
\right].
\tag{24}
\]

These relations lead naturally to the quadratic wave equation.

But consider the following modification:
\[
i\hbar\partial_s\psi
=
\left[
\frac{1}{2m}\hat\pi_\mu\hat\pi^\mu
+W(x,s)
\right]\psi,
\tag{25}
\]

where \(W\) is any real-valued scalar potential.

The additional term does not change the continuity equation or the probability current (24).

Thus probability conservation alone cannot distinguish between these generators.

This is not a minor technical objection.

It means that the probability argument does not uniquely determine the evolution generator.

There are further possibilities involving different differential operators and corresponding probability currents.
What does this imply?
I would reformulate Fanchi's conclusion as follows:
- The quadratic generator is naturally compatible with a positive-definite spacetime probability density.
- It has a simple classical Hamiltonian counterpart.
- It produces a conventional second-order covariant wave equation.
- It is compatible with electromagnetic gauge coupling.
- It has a straightforward nonrelativistic limit.
These are substantial reasons for choosing it.

But the statement that it follows uniquely from probability theory requires additional assumptions about the current, the differential order of the equation, and the classical correspondence.

This is where I would concentrate a critical examination of the 1994 paper.

8. Another important issue: the relationship with Klein–Gordon theory
There is a particularly illuminating way to understand the quadratic generator.

Take the free equation
\[
i\hbar\partial_s\psi
=
\frac{\hat p_\mu\hat p^\mu}{2m}\psi.
\tag{26}
\]

Look for solutions of the form
\[
\psi(x,s)=\phi(x)e^{-i\kappa s/\hbar}.
\]

Substitution gives
\[
\frac{\hat p_\mu\hat p^\mu}{2m}\phi
=
\kappa\phi.
\]

Consequently,
\[
\boxed{
\left(
\hat p_\mu\hat p^\mu-2m\kappa
\right)\phi=0.
}
\tag{27}
\]

If we choose
\[
2m\kappa=m_0^2c^2,
\]

we recover the Klein–Gordon equation for mass \(m_0\):
\[
\boxed{
\left(
\hat p_\mu\hat p^\mu-m_0^2c^2
\right)\phi=0.
}
\tag{28}
\]

Here I use the metric convention
\[
\eta_{\mu\nu}=\operatorname{diag}(+,-,-,-).
\]

This relationship is important.

The Klein–Gordon equation appears as a fixed-eigenvalue sector of the parametrized theory.

The parametrized theory, however, permits a superposition of different eigenvalues:
\[
\psi(x,s)=
\int d\kappa\,
a(\kappa)\phi_\kappa(x)
e^{-i\kappa s/\hbar}.
\]

Therefore, the invariant mass need not be fixed throughout the evolution.

This is precisely the off-shell feature that distinguishes parametrized relativistic dynamics from conventional single-mass relativistic quantum mechanics.

9. My critical assessment
I would separate the achievements and unresolved issues as follows.
Issue	Assessment
Positive spacetime probability density	Well-defined mathematical proposal
Conservation of probability	Consistent with the parametrized equation
Quadratic generator	Natural and physically motivated
Unique derivation from probability	Not established by continuity alone
Classical correspondence	Particularly attractive
Off-shell mass spectrum	Important additional physical structure
Experimental confirmation	Still an open question

There is also a question about the Hilbert-space interpretation.

The proposed scalar product is
\[
\langle\phi|\psi\rangle
=
\int_{\mathbb R^4}
\phi^*(x)\psi(x)\,d^4x.
\tag{29}
\]

This is positive definite, unlike the familiar Klein–Gordon inner product.

But the operator
\[
\hat p_\mu\hat p^\mu
\]

has a spectrum unbounded both above and below. This does not prevent unitary evolution in \(s\), provided the generator is self-adjoint, but it raises further questions about stability and the physical interpretation of the spectrum.

I would not regard the positive-definite norm as resolving all probabilistic problems automatically.

10. Why this matters for our Nelson–SHP project
Here I see a particularly interesting possibility.

Fanchi's construction gives us a positive probability distribution
\[
\rho(x,s)=|\psi(x,s)|^2
\]

and a current
\[
J^\mu=\rho V^\mu.
\]

Nelson's mechanics, on the other hand, starts with forward and backward stochastic processes.

In nonrelativistic notation, these give two mean derivatives,
\[
D_+,\qquad D_-,
\]

and two associated velocities,
\[
b_+,\qquad b_-.
\]

Their average and difference define the current and osmotic velocities:
\[
v=\frac{b_++b_-}{2},
\qquad
u=\frac{b_+-b_-}{2}.
\]

The central Nelson relation is
\[
v=\frac{\nabla S}{m}.
\]

Fanchi's relativistic expression
\[
V^\mu=
\frac{1}{m}
\left(
\partial^\mu S-\frac ec A^\mu
\right)
\]

looks like a natural starting point for a covariant generalization.

But there is a serious issue: a relativistic diffusion must preserve the appropriate causal and geometric structure. One cannot simply replace the spatial Laplacian in Nelson's equations by the d'Alembertian and expect a valid stochastic process.

I would therefore proceed in the following order.
A possible research programme
1. Fanchi's probabilistic foundation
   
   Establish the spacetime probability measure and its evolution in \(s\).
2. Relativistic stochastic kinematics
   
   Construct forward and backward processes compatible with the invariant parameter.
3. Mass-shell dynamics
   
   Determine whether the stochastic process should evolve on the mass shell or in an extended off-shell phase space.
4. Quantum dynamics
   
   Derive the parametrized wave equation rather than assuming it.
5. Spin and gauge fields
   
   Introduce spinor bundles and stochastic parallel transport.

11. What I would investigate next
There is one question that I think should be settled before proceeding:

Can Fanchi's quadratic generator be derived from a more general probabilistic principle, rather than being selected by assuming a particular form of probability current?

One possible route is to formulate the problem in terms of:
1. A positive probability density on spacetime.
2. A local conservation law in \(s\).
3. Lorentz covariance.
4. Gauge invariance.
5. A suitable classical limit.
6. A specified differential order and locality requirement.
Then determine which generators are actually allowed.

This would be a useful foundational result in its own right. It would also tell us exactly which assumptions a subsequent relativistic Nelson theory must inherit.

My recommendation is to pursue this question first. It prevents us from building a stochastic extension on a generator whose purported uniqueness has not yet been established.
Toward a probabilistic foundation of parametrized relativistic dynamics
Let us pursue the question we identified: under what assumptions does the quadratic evolution generator follow from probability theory, Lorentz covariance and classical mechanics?

I think we can make substantial progress without committing ourselves to Fanchi's particular argument.

The key result will be:
\[
\boxed{
\text{The quadratic generator is essentially selected by locality,}
\atop
\text{second-order dynamics, covariance and classical correspondence.}
}
\]

Probability conservation is necessary, but insufficient by itself.

I will first consider a spinless particle in flat spacetime, using your preferred convention
\[
\eta_{\mu\nu}=\operatorname{diag}(+1,-1,-1,-1).
\]

1. General formulation
Let the wavefunction be
\[
\psi=\psi(x^\mu,s),
\]

with probability density
\[
\rho(x,s)=|\psi(x,s)|^2.
\]

We require
\[
\int_{\mathbb R^4}\rho(x,s)\,d^4x=1.
\tag{1}
\]

The evolution equation is
\[
i\hbar\partial_s\psi=K\psi.
\tag{2}
\]

For conservation of the spacetime probability, we require
\[
\frac{\partial\rho}{\partial s}
+\partial_\mu J^\mu=0.
\tag{3}
\]

Integration over spacetime gives
\[
\frac{d}{ds}\int\rho\,d^4x
=
-\int\partial_\mu J^\mu\,d^4x=0,
\]

assuming vanishing flux at infinity.

Notice that this construction does not require a positive-definite spacetime metric. Positivity comes from \(|\psi|^2\), not from the Minkowski metric.

This is an important advantage over the conventional Klein–Gordon probability density.
2. Classifying possible evolution generators
Suppose we impose the following requirements:
1. \(K\) is a Lorentz scalar.
2. \(K\) is local in spacetime.
3. \(K\) is a differential operator of at most second order in \(x^\mu\).
4. \(K\) is Hermitian with respect to the spacetime scalar product.
5. \(K\) is invariant under electromagnetic gauge transformations.
6. Its classical limit describes a relativistic particle.
Introduce the kinetic momentum operator
\[
\hat\pi_\mu
=
-i\hbar\partial_\mu-\frac qc A_\mu.
\tag{4}
\]

Here \(q\) denotes the electric charge.

The most general scalar second-order operator with the appropriate principal part has the schematic form
\[
\boxed{
K=
\frac{1}{2M}\hat\pi_\mu\hat\pi^\mu
+W(x,s).
}
\tag{5}
\]

Here \(M\) is a parameter with dimensions of mass, and \(W\) is a real scalar interaction.

There can also be ordering terms if the coefficients depend on position, and additional curvature terms in curved spacetime.

Why does this form emerge?

The principal symbol of a second-order Lorentz-invariant scalar differential operator must be proportional to
\[
\eta^{\mu\nu}p_\mu p_\nu.
\]

There is no other independent quadratic Lorentz scalar constructed from a single four-momentum.

Gauge invariance then requires replacing canonical momentum by kinetic momentum:
\[
p_\mu\longrightarrow p_\mu-\frac qc A_\mu.
\]

Finally, classical correspondence fixes the coefficient of the quadratic term to \(1/(2M)\).

This is already a rather strong argument for Fanchi's choice.

But notice what has happened: we have used locality, differential order and classical correspondence in addition to probability conservation.
3. Probability conservation follows automatically
Consider
\[
i\hbar\partial_s\psi
=
\left(
\frac{\hat\pi_\mu\hat\pi^\mu}{2M}+W
\right)\psi.
\tag{6}
\]

For real \(W\), multiplication by \(W\) makes no contribution to the difference between the equation and its complex conjugate.

Define
\[
J^\mu=
\frac{1}{2M}
\left[
\psi^*\hat\pi^\mu\psi+
(\hat\pi^\mu\psi)^*\psi
\right].
\tag{7}
\]

A direct calculation gives
\[
\boxed{
\partial_s|\psi|^2+\partial_\mu J^\mu=0.
}
\tag{8}
\]

Thus the quadratic generator has precisely the desired probability conservation law.

But the converse does not hold: equation (8) does not uniquely determine equation (6).

This is the distinction I would insist on when evaluating Fanchi's foundational argument.
4. The classical limit provides an additional selection principle
Now let us use the polar representation
\[
\psi=R\exp(iS/\hbar),
\qquad \rho=R^2.
\]

Substitution into equation (6), followed by separation into real and imaginary parts, gives two equations.

The imaginary part is the continuity equation:
\[
\partial_s\rho+
\partial_\mu\left[
\frac{\rho}{M}
\left(\partial^\mu S-\frac qc A^\mu\right)
\right]=0.
\tag{9}
\]

The real part is
\[
\boxed{
\partial_s S+
\frac{1}{2M}
\left(\partial_\mu S-\frac qc A_\mu\right)
\left(\partial^\mu S-\frac qc A^\mu\right)
+W
-\frac{\hbar^2}{2M}
\frac{\Box R}{R}
=0.
}
\tag{10}
\]

This is a parametrized relativistic Hamilton–Jacobi equation with an additional quantum potential,
\[
Q=-\frac{\hbar^2}{2M}\frac{\Box R}{R}.
\tag{11}
\]

In the classical limit, neglecting \(Q\), we recover
\[
\partial_s S+K_{\rm cl}(x,\partial S,s)=0,
\tag{12}
\]

where
\[
K_{\rm cl}
=
\frac{1}{2M}
\left(p_\mu-\frac qc A_\mu\right)
\left(p^\mu-\frac qc A^\mu\right)
+W.
\tag{13}
\]

This is precisely the Hamilton–Jacobi equation associated with the classical parametrized relativistic Hamiltonian.

The corresponding Hamilton equations are
\[
\frac{dx^\mu}{ds}
=
\frac{1}{M}
\left(p^\mu-\frac qc A^\mu\right),
\tag{14}
\]

and
\[
\frac{dp_\mu}{ds}
=
-\frac{\partial K_{\rm cl}}{\partial x^\mu}.
\tag{15}
\]

Thus the quadratic generator is not an arbitrary relativistic generalization of Schrödinger evolution. It is the direct quantization of a familiar classical Hamiltonian.

This is, in my view, a stronger justification than probability conservation alone.
5. What is actually unique?
We can now formulate a more precise result.
Conditional uniqueness statement

Consider a spinless relativistic particle with:
- A scalar wavefunction on Minkowski spacetime.
- First-order evolution in an invariant parameter \(s\).
- A local, second-order differential generator.
- Lorentz covariance and electromagnetic gauge invariance.
- A positive-definite spacetime probability norm.
- Classical correspondence with a quadratic relativistic Hamiltonian.
Then, up to scalar interaction terms and choices of operator ordering, the generator is
\[
K=\frac{1}{2M}\hat\pi_\mu\hat\pi^\mu+W.
\]

This is a conditional uniqueness result, not a uniqueness theorem based on probability conservation alone.

There are several ways to go beyond this result.
5.1 Higher-order generators
We could consider
\[
K=f(\hat\pi_\mu\hat\pi^\mu).
\]

For example,
\[
K=
\frac{\hat\pi^2}{2M}
+\alpha\frac{(\hat\pi^2)^2}{M^3c^2}.
\]

This remains Lorentz and gauge invariant, and can be Hermitian under suitable domain conditions.

It also generates unitary evolution in \(s\), but it introduces higher-order spacetime derivatives.

Its classical limit is no longer the ordinary quadratic relativistic Hamiltonian.

Consequently, classical correspondence and the requirement of second-order spacetime dynamics exclude it.
5.2 Square-root generators
Another possibility is
\[
K=\sqrt{\hat\pi^2}.
\]

This is generally nonlocal in spacetime and requires a careful specification of its spectral domain.

It is not excluded by probability conservation in principle, but it does not satisfy our locality requirement.
5.3 Additional scalar interactions
As we have already seen,
\[
K=\frac{\hat\pi^2}{2M}+W(x,s)
\]

has exactly the same continuity equation for any real \(W\).

This means that probability conservation cannot determine the interaction potential.

The potential must be specified by additional physical principles.

6. An overlooked extension: curved spacetime
This is particularly relevant to the SHP programme.

In curved spacetime, the natural generalization is
\[
\boxed{
K=
-\frac{\hbar^2}{2M}
\frac{1}{\sqrt{|g|}}
D_\mu
\left[
\sqrt{|g|}\,g^{\mu\nu}D_\nu
\right]
+W.
}
\tag{16}
\]

Here \(D_\mu\) is the electromagnetic covariant derivative.

The probability norm becomes
\[
\int |\psi|^2\sqrt{|g|}\,d^4x=1.
\tag{17}
\]

The corresponding conservation law is
\[
\partial_s\rho+
\nabla_\mu J^\mu=0,
\tag{18}
\]

where \(\rho=|\psi|^2\) and
\[
J^\mu=
\frac{1}{2M}
\left[
\psi^*\hat\pi^\mu\psi+
(\hat\pi^\mu\psi)^*\psi
\right].
\]

There is an additional freedom:
\[
K\longrightarrow K+\frac{\hbar^2\xi}{2M}R,
\tag{19}
\]

where \(R\) is the scalar curvature and \(\xi\) is dimensionless.

This term is:
- Lorentz covariant.
- Gauge invariant.
- Local.
- Hermitian.
- Compatible with probability conservation.
Yet it cannot be determined by the preceding argument.

This is a useful illustration of why the quadratic kinetic term may be selected without the entire evolution generator being uniquely fixed.
7. What happens when we introduce spin?
Here the problem becomes considerably richer.

Suppose that instead of a scalar wavefunction we introduce a spinor
\[
\psi_\alpha(x,s).
\]

There are at least two possible constructions.

Option A: retain the quadratic generator

We can use
\[
K=\frac{1}{2M}\hat\pi_\mu\hat\pi^\mu.
\]

But this operator acts identically on all spinor components. Merely making the wavefunction a spinor does not introduce spin dynamics.

Option B: construct the generator from the Dirac operator

Define
\[
\mathcal D=\gamma^\mu\hat\pi_\mu.
\]

Then consider
\[
K=\frac{1}{2M}\mathcal D^2.
\tag{20}
\]

This is much more interesting.

Because the kinetic momenta do not commute in an electromagnetic field,
\[
[\hat\pi_\mu,\hat\pi_\nu]
=
i\hbar\frac qc F_{\mu\nu},
\]

the square of the Dirac operator contains a spin-field interaction.

Indeed,
\[
\mathcal D^2
=
\hat\pi_\mu\hat\pi^\mu
+
\frac{i\hbar q}{4c}
[\gamma^\mu,\gamma^\nu]F_{\mu\nu}.
\tag{21}
\]

Writing
\[
\sigma^{\mu\nu}
=
\frac{i}{2}[\gamma^\mu,\gamma^\nu],
\]

we obtain
\[
\boxed{
K=
\frac{1}{2M}
\left(
\hat\pi^2+
\frac{\hbar q}{2c}\sigma^{\mu\nu}F_{\mu\nu}
\right).
}
\tag{22}
\]

The overall sign of the spin term depends on the convention for charge, covariant momentum and the definition of the Dirac operator; equation (22) uses the conventions specified above.

This is an important observation: the Pauli interaction can emerge from the noncommutativity of covariant momentum components.

It is not necessary to insert a magnetic-moment interaction by hand.

However, this construction is not yet a complete stochastic theory of spin. We still need to establish the appropriate spinor probability current, the physical meaning of the spinor components, and the relationship between the parametrized equation and the first-order Dirac equation.

For a discussion of the latter, see the classic work of Feynman and Gell-Mann on the second-order formulation of fermion dynamics.
8. A significant distinction: the mass parameter versus physical mass
I want to emphasize one point that can easily become confusing in Fanchi's formulation.

In
\[
K=\frac{\hat\pi^2}{2M},
\]

the parameter \(M\) determines the coefficient of the kinetic term.

It need not be identical to the invariant mass of every state.

For a free particle, consider an eigenstate satisfying
\[
\hat\pi^2\phi_\kappa=2M\kappa\phi_\kappa.
\]

Its evolution is
\[
\psi_\kappa(x,s)
=
\phi_\kappa(x)e^{-i\kappa s/\hbar}.
\]

If the state is timelike, we can identify its invariant mass \(m_{\rm inv}\) through
\[
2M\kappa=m_{\rm inv}^2c^2.
\]

Thus
\[
m_{\rm inv}^2=\frac{2M\kappa}{c^2}.
\]

The theory can therefore contain a distribution of invariant masses.

This is an essential aspect of off-shell dynamics.

It also raises a question that deserves more attention than it usually receives: what physical mechanism determines the distribution of mass eigenvalues, and why should a stable particle have a sharply defined mass?

9. Where I think we should go next
We now have a reasonably clear hierarchy:
Principle	What it determines
Probability conservation	A continuity equation
Positive spacetime norm	A probabilistic interpretation
Lorentz covariance	The allowed tensor structures
Gauge invariance	The electromagnetic covariant momentum
Locality and second-order dynamics	Quadratic principal operator
Classical correspondence	Its kinetic coefficient
Additional physics	Scalar potentials, curvature coupling, spin interactions

This suggests that Fanchi's quadratic generator is well motivated, but its selection is conditional on a collection of physical assumptions.

I would now take the next step in a direction that is particularly relevant to our original project.
The next fundamental problem: relativistic Nelson mechanics

Can we construct a stochastic process whose probability density satisfies
\[
\partial_s\rho+\partial_\mu(\rho v^\mu)=0
\]

and whose stochastic dynamics produces
\[
i\hbar\partial_s\psi
=
\frac{\hat\pi^2}{2M}\psi?
\]

More specifically:
1. What is the appropriate relativistic diffusion process?
2. What replaces the ordinary spatial osmotic velocity?
3. How do we avoid superluminal stochastic trajectories?
4. Can the mass-shell condition emerge statistically rather than being imposed?
5. Can spin arise from stochastic parallel transport?

My preference would be to investigate the third and fourth questions together.

There is a fundamental tension here: ordinary Nelson diffusion is Brownian and has nondifferentiable paths, whereas relativistic particle trajectories are usually required to be timelike. A diffusion directly in spacetime does not automatically preserve that requirement.

One possible way forward is to formulate the stochastic process in an extended phase space, with position and four-momentum as stochastic variables, and constrain the momentum to the future mass hyperboloid when appropriate.

But this introduces a new question: should we abandon Fanchi's off-shell dynamics, or should we allow stochastic excursions away from the mass shell?

I suspect the latter is more compatible with the original SHP–Fanchi motivation.

That is where I would place the next serious investigation.
Relativistic Nelson mechanics in extended phase space: an off-shell construction
Let us develop this as an actual dynamical model, rather than simply replace the Laplacian in Nelson's equations by the d'Alembertian.

I propose the following starting point:

The stochastic process should live in relativistic phase space, not directly in spacetime.

Its basic variables are
\[
Z_s=(X_s^\mu,P_s^\mu),
\]

where \(s\) is the invariant evolution parameter.

This has three advantages:
1. The position evolves continuously and can have a timelike tangent.
2. Randomness acts on momentum rather than directly on spacetime position.
3. We can retain the off-shell structure of Fanchi and SHP.
There is a price: we must construct a suitable stochastic dynamics for momentum, and we must distinguish the instantaneous mass shell from the statistical distribution of invariant masses.

I will first construct the spinless theory in flat spacetime, with electromagnetic interactions. Then we can examine its relation to Nelson's equations and identify what remains to be proved.
1. The classical phase-space structure
Use
\[
\eta_{\mu\nu}=\operatorname{diag}(+1,-1,-1,-1).
\]

Define the kinetic momentum
\[
\pi^\mu=P^\mu-\frac qc A^\mu(X).
\]

The natural classical Hamiltonian is
\[
K(X,P)=\frac{1}{2M}\pi_\mu\pi^\mu+q\phi_s(X),
\]

where \(\phi_s\) denotes a possible additional scalar interaction (not the electromagnetic scalar potential, which is already contained in \(A_\mu\)).

For simplicity, initially put \(\phi_s=0\).

Hamilton's equations are
\[
\frac{dX^\mu}{ds}=\frac{\pi^\mu}{M},
\tag{1}
\]

\[
\frac{dP_\mu}{ds}
=
\frac q{Mc}\pi^\nu\partial_\mu A_\nu.
\tag{2}
\]

Combining them gives
\[
\boxed{
M\frac{d^2X^\mu}{ds^2}
=
\frac qc F^\mu{}_\nu
\frac{dX^\nu}{ds}.
}
\tag{3}
\]

This is the Lorentz-force equation parametrized by \(s\).

The important observation is that equation (3) does not require
\[
\pi_\mu\pi^\mu=M^2c^2.
\]

Indeed,
\[
\frac{d}{ds}(\pi_\mu\pi^\mu)
=
\frac{2q}{c}\pi_\mu F^\mu{}_\nu\pi^\nu=0.
\tag{4}
\]

Thus the invariant mass is conserved along a free or electromagnetically interacting trajectory, but its value need not equal the reference mass \(M\).

This is precisely the kind of off-shell dynamics we want.
2. Why stochasticity should act on momentum
Suppose we introduce ordinary Brownian noise directly into position:
\[
dX^\mu=b^\mu ds+\sqrt{2D}\,dW^\mu.
\]

This is problematic.

The quadratic variation is
\[
d[X^\mu,X^\nu]_s
=
2D\eta^{\mu\nu}ds
\]

if one tries to use a Lorentz-invariant Wiener process. But the Minkowski metric is indefinite, and there is no ordinary positive probability measure associated with such a covariance matrix.

Furthermore, Brownian paths are nondifferentiable. They do not have an ordinary four-velocity, and their tangents cannot be required to remain timelike in the usual sense.

The natural alternative is
\[
\boxed{
dX^\mu_s=\frac{P^\mu_s-qA^\mu(X_s)/c}{M}\,ds.
}
\tag{5}
\]

Position has finite variation, while momentum is stochastic.

This is similar in spirit to kinetic Brownian motion, where the position is obtained by integrating a stochastic velocity.

It also preserves the possibility of defining a timelike trajectory, provided
\[
\pi_\mu\pi^\mu>0,\qquad \pi^0>0.
\]

Notice that this is not yet a Nelson diffusion. It is a kinetic diffusion whose stochasticity acts on momentum.

That distinction will matter later.
3. The correct relativistic momentum diffusion
There is a subtlety in the proposal above.

Adding an unconstrained four-dimensional Wiener noise to \(P^\mu\) does not guarantee that \(P^\mu\) remains future timelike. Consequently, even if the position equation is deterministic in differential form, the resulting trajectories can become spacelike.

We need a more careful construction.

Introduce a decomposition
\[
\boxed{
P^\mu_s=\mu_s c\,U^\mu_s,
\qquad
U_\mu U^\mu=1,
\qquad U^0>0.
}
\tag{6}
\]

Here:
- \(U^\mu\) is a dimensionless future-directed unit four-velocity.
- \(\mu_s>0\) is the instantaneous invariant mass.
- \(P^\mu\) is the kinetic four-momentum.
The unit vectors \(U^\mu\) belong to the future unit hyperboloid
\[
H^3=\{U:U_\mu U^\mu=1,\ U^0>0\}.
\]

This is a Riemannian manifold, despite being embedded in Minkowski spacetime.

Its induced metric is positive definite.

That observation is crucial: we can have ordinary positive-probability Brownian diffusion on the velocity hyperboloid without introducing a Wiener process with an indefinite Minkowski covariance.







commons.wikimedia.org


3.1 Diffusion on the hyperboloid
Let \(W_s\) be a three-dimensional Brownian motion in the local tangent space of \(H^3\).

The geometrically invariant diffusion has generator
\[
\mathcal L_U
=
\frac{\nu}{2}\Delta_{H^3},
\tag{7}
\]

where \(\nu\) is a diffusion coefficient and \(\Delta_{H^3}\) is the Laplace–Beltrami operator.

The corresponding stochastic process can be written geometrically as
\[
dU_s=\text{parallel-transported noise}
+\text{drift tangent to }H^3.
\]

The tangency condition is
\[
U_\mu\,dU^\mu
+
\frac12 d[U_\mu,U^\mu]=0.
\tag{8}
\]

This is the stochastic counterpart of differentiating \(U^2=1\).

The resulting process stays on the future unit hyperboloid.
3.2 Allowing the mass to fluctuate
Now introduce a positive stochastic mass:
\[
\mu_s=\mu_*e^{Y_s}.
\]

For example,
\[
dY_s=a(Y_s,U_s,X_s)\,ds+\sigma\,dB_s.
\tag{9}
\]

This guarantees \(\mu_s>0\).

The full momentum process is then
\[
P^\mu_s=\mu_*c\,e^{Y_s}U^\mu_s.
\tag{10}
\]

It follows that
\[
P_\mu P^\mu=\mu_s^2c^2>0
\]

at every instant.

Thus the particle remains timelike, while its invariant mass can fluctuate.

This is a useful compromise between two requirements:
- Preserve causal particle trajectories.
- Retain the off-shell mass degree of freedom.
There is one important qualification. The mass diffusion in (9) is an additional physical postulate. It does not follow from ordinary four-dimensional electromagnetic coupling. In the simplest SHP electromagnetic dynamics, the mass-shell function is conserved along a trajectory. Genuine mass-changing dynamics requires an additional mechanism, for example a fifth gauge interaction or an explicitly specified stochastic force.
4. The resulting position process
Define
\[
\frac{dX^\mu_s}{ds}
=
\frac{P^\mu_s}{M}.
\]

Using (10),
\[
\boxed{
dX^\mu_s=
\frac{\mu_*c}{M}e^{Y_s}U^\mu_s\,ds.
}
\tag{11}
\]

Since
\[
U^2=1,\qquad U^0>0,
\]

we have
\[
\left(\frac{dX}{ds}\right)^2
=
\frac{\mu_s^2c^2}{M^2}>0.
\]

Therefore the position trajectories are future-directed and timelike.

They are continuously differentiable if the momentum process has continuous paths.

This solves the causal-trajectory problem that occurs when one simply writes a four-dimensional Brownian motion for \(X^\mu\).

But it introduces another distinction: this is a relativistic kinetic diffusion, not yet a relativistic Nelson theory.

5. The phase-space probability equation
Let
\[
f(x,\mu,u,s)
\]

be the probability density on the extended phase space
\[
\mathcal P=\mathbb R^{1,3}\times\mathbb R_+\times H^3.
\]

The natural measure is
\[
d\Gamma=d^4x\,d\mu\,d\operatorname{vol}_{H^3}.
\]

For simplicity, consider independent mass and velocity diffusions, with no external fields initially.

The forward Fokker–Planck equation has the schematic form
\[
\boxed{
\begin{aligned}
\partial_s f
&+
\frac{\mu c}{M}u^\mu\partial_\mu f\\
&=
-\partial_\mu^{(m)}(a_m f)
+\frac12\partial_\mu^{(m)2}(D_m f)
\\
&\quad
-\operatorname{div}_{H^3}(b f)
+\frac{\nu}{2}\Delta_{H^3}f.
\end{aligned}
}
\tag{12}
\]

Here \(a_m\) and \(D_m\) denote the mass drift and diffusion coefficients (with the mass-coordinate derivative understood), while \(b\) is a tangent drift on the hyperboloid.

The equation describes a normalized probability distribution:
\[
\int_{\mathcal P} f\,d\Gamma=1.
\]

Integrating over \(\mu\) and \(u\), we obtain a spacetime density
\[
\rho(x,s)=
\int_0^\infty d\mu
\int_{H^3} f(x,\mu,u,s)\,d\operatorname{vol}_{H^3}.
\]

Its current is
\[
J^\mu(x,s)=
\int_0^\infty d\mu
\int_{H^3}
\frac{\mu c}{M}u^\mu f\,
d\operatorname{vol}_{H^3}.
\]

Consequently,
\[
\boxed{
\partial_s\rho+\partial_\mu J^\mu=0.
}
\tag{13}
\]

This is exactly the type of probability conservation equation that we need for a parametrized relativistic theory.

However, it does not yet imply a Schrödinger equation. Equation (12) is a real kinetic equation, whereas Fanchi's equation is a complex, unitary evolution equation.

We need an additional step.
6. How do we recover Nelson's forward and backward structure?
Now comes the central step.

For a Markov diffusion on a Riemannian configuration manifold, the forward and backward mean derivatives are defined by
\[
D_+F(Z_s)
=
\lim_{h\downarrow0}
E\left[
\frac{F(Z_{s+h})-F(Z_s)}{h}
\mid Z_s
\right],
\]

\[
D_-F(Z_s)
=
\lim_{h\downarrow0}
E\left[
\frac{F(Z_s)-F(Z_{s-h})}{h}
\mid Z_s
\right].
\]

Here \(Z_s=(X_s,Y_s,U_s)\).

For the stochastic variables \(Y,U\), these derivatives differ because of diffusion.

For the position variable \(X\), however, we have
\[
dX^\mu=V^\mu(Z_s)\,ds.
\]

Consequently,
\[
D_+X^\mu=D_-X^\mu=V^\mu.
\tag{14}
\]

This is a significant departure from ordinary Nelson mechanics.

There is no position osmotic velocity:
\[
u_X^\mu
=
\frac12(D_+X^\mu-D_-X^\mu)=0.
\]

The osmotic dynamics resides in momentum space.

This suggests a different formulation:

Nelson's forward/backward structure should be applied to the extended phase-space variables, not to spacetime position alone.
6.1 Forward and backward momentum drifts
Let \(Q^a\) denote coordinates on the internal stochastic manifold
\[
\mathcal N=\mathbb R\times H^3,
\]

where the first coordinate is \(Y=\log(\mu/\mu_*)\).

Let \(h_{ab}\) be its Riemannian metric. For simplicity, take
\[
h_{ab}dQ^a dQ^b
=
dY^2+\ell^2 ds_{H^3}^2,
\]

where \(\ell\) is a constant scale.

Consider the diffusion generator
\[
L_+
=
V^\mu(Q)\partial_\mu
+b_+^a\nabla_a
+\frac{\nu}{2}\Delta_{\mathcal N}.
\tag{15}
\]

Here \(b_+^a\) is the forward drift.

If \(f\) is the probability density relative to the invariant measure on \(\mathcal N\), the backward drift is
\[
\boxed{
b_-^a
=
b_+^a-\nu\nabla^a\ln f.
}
\tag{16}
\]

This is the manifold version of Nelson's relation between forward and backward drifts.

We can define
\[
v^a=\frac12(b_+^a+b_-^a),
\]

\[
u^a=\frac12(b_+^a-b_-^a).
\]

Then
\[
\boxed{
u^a=\frac{\nu}{2}\nabla^a\ln f.
}
\tag{17}
\]

Thus the osmotic velocity is now a vector field in the mass-velocity sector.

The corresponding mean acceleration is
\[
\mathcal A^a
=
\frac12(D_+D_-+D_-D_+)Q^a.
\tag{18}
\]

In a curved stochastic manifold, this expression must be interpreted using covariant stochastic derivatives and the associated curvature terms.
6.2 The stochastic Newton equation
The natural relativistic analogue of Nelson's dynamical equation would be
\[
\boxed{
M\mathcal A^\mu
=
\frac qc F^\mu{}_\nu V^\nu.
}
\tag{19}
\]

But we must be careful: this is not yet a complete equation for the internal variables \(Y,U\).

We need to supplement it with equations governing mass diffusion and hyperbolic momentum diffusion.

A more complete formulation is variational:
\[
\delta E\left[
\int ds\,
L_{\rm classical}(X,Q,D_+Z,D_-Z)
\right]=0,
\tag{20}
\]

where the stochastic action is the expectation of the classical action, and the variations are restricted to admissible processes on the phase-space manifold.

This is closely related to the approach of Marra and Serva (1990), who constructed relativistic stochastic variational principles and derived stochastic Lagrange equations. 
Numdam


7. An important historical correction: this is not an entirely unexplored direction
There is relevant earlier literature that we should incorporate rather than reinvent.

In particular:

Marra and Serva (1990)

Variational Principles for a Relativistic Stochastic Mechanics, Annales de l'Institut Henri Poincaré A 53, 97–108.

They investigate relativistic covariant stochastic processes, mean classical actions, forward/backward constructions and the recovery of Klein–Gordon dynamics. Their introduction also discusses the earlier work of Guerra–Ruggiero and Dohrn–Guerra. 
Numdam


Relativistic random flights (1992)

Relativistic Problem of Random Flights and Nelson's Stochastic Mechanics, Physics Letters A 164, 6–16.

This work connects phase-space random flights with the diffusion limit and Klein–Gordon dynamics. 
ScienceDirect


Z. Haba (2009)

Relativistic Diffusion, Physical Review E 79, 021128.

This is particularly relevant to our hyperbolic momentum-space construction. It develops proper-time relativistic diffusion, electromagnetic interactions and kinetic equations on the relativistic momentum geometry. 
Phys. Rev. E


These papers establish that relativistic stochastic mechanics has a history considerably richer than is sometimes acknowledged.

Our proposed construction differs in its emphasis on retaining an independent stochastic mass variable while keeping the position trajectories timelike.

That combination needs a careful comparison with the existing approaches.
8. Can this construction produce Fanchi's wave equation?
This is the decisive test.

Recall Fanchi's equation:
\[
i\hbar\partial_s\psi
=
\frac{1}{2M}\hat\pi_\mu\hat\pi^\mu\psi.
\tag{21}
\]

Its polar decomposition produces the continuity equation and the quantum Hamilton–Jacobi equation.

We would like our stochastic phase-space theory to produce the same pair of equations.

However, there is an obstacle.

Our phase-space diffusion produces a kinetic equation of Kramers type:
\[
\partial_s f+
\frac{\mu c}{M}u^\mu\partial_\mu f
=
\mathcal L_{\rm internal}^*f.
\tag{22}
\]

After integration over momentum, we obtain
\[
\partial_s\rho+\partial_\mu J^\mu=0.
\]

But the equation for \(J^\mu\) involves higher momentum moments. Consequently, the equation for \(\rho\) does not close automatically.

This is the familiar moment-closure problem of kinetic theory.

In particular, we cannot simply identify
\[
f(x,\mu,u,s)
\longrightarrow |\psi(x,s)|^2
\]

and claim to have derived Fanchi's equation.

We need an additional dynamical principle.
A possible route: stochastic Hamilton–Jacobi theory
Introduce a phase function \(S(x,s)\) and require that the current velocity satisfy
\[
J^\mu
=
\frac{\rho}{M}
\left(
\partial^\mu S-\frac qc A^\mu
\right).
\tag{23}
\]

Then the continuity equation becomes
\[
\partial_s\rho+
\partial_\mu
\left[
\frac{\rho}{M}
\left(\partial^\mu S-\frac qc A^\mu\right)
\right]=0.
\tag{24}
\]

We additionally seek a stochastic variational principle yielding
\[
\partial_s S+
\frac{1}{2M}
\left(\partial_\mu S-\frac qc A_\mu\right)
\left(\partial^\mu S-\frac qc A^\mu\right)
+Q=0.
\tag{25}
\]

If the stochastic variation produces
\[
Q=-\frac{\hbar^2}{2M}
\frac{\Box\sqrt{\rho}}{\sqrt{\rho}},
\tag{26}
\]

then equations (24) and (25) combine into Fanchi's wave equation.

But equation (26) must be derived from the stochastic action. We cannot insert it simply because it gives the desired answer.

There is also a conceptual question: our stochastic process lives on \(\mathbb R^{1,3}\times\mathbb R_+\times H^3\), whereas the quantum potential in (26) involves the spacetime d'Alembertian. We must show how eliminating the internal variables produces precisely that operator.

This is the most difficult part of the proposal.

9. Three possible formulations
I would now distinguish three approaches rather than force everything into one construction.
Approach	Position trajectories	Mass shell	Main challenge
Spacetime diffusion	Brownian, nondifferentiable	Not automatically preserved	Causality and positive diffusion
Fixed-mass hyperbolic diffusion	Timelike, differentiable	Fixed	Recovering quantum wave dynamics
Extended phase-space diffusion	Timelike, differentiable	Dynamical mass	Deriving the quantum potential

The third approach is the one I would investigate for an SHP–Fanchi synthesis.

But there is an important alternative worth keeping in mind: the relativistic diffusion literature often uses a fixed mass hyperboloid, while the SHP programme naturally allows off-shell motion. These need not be mutually exclusive, but their combination requires a genuinely new dynamical construction.
10. What about curved spacetime?
The extension is conceptually straightforward.

Replace the future unit hyperboloid by the future unit tangent bundle
\[
UH(M)=
\{(x,u):g_{\mu\nu}u^\mu u^\nu=1,\ u\text{ future directed}\}.
\]

The position equation becomes
\[
\frac{dX^\mu}{ds}
=
\frac{\mu c}{M}U^\mu.
\]

The velocity evolves through a combination of:
- Geodesic transport.
- Electromagnetic Lorentz force.
- Stochastic diffusion tangent to the velocity hyperboloid.
- Possible mass fluctuations.
The natural geometric tools are the Levi-Civita connection, horizontal and vertical vector fields on the tangent bundle, and the associated horizontal and vertical Laplacians.

This is closely related to the geometric framework developed by Franchi and Le Jan, although our proposed additional mass variable changes the construction.

For the relevant background, see:
- Franchi and Le Jan, Relativistic Diffusions.
- Haba, Relativistic Diffusion with Friction on a Pseudo-Riemannian Manifold (2010).
- Herrmann, Diffusion in General Relativity (2010).

11. My assessment and the next concrete calculation
I think we have identified a worthwhile problem, but I would not yet call the proposed theory a derivation of relativistic Nelson mechanics.

The most important results so far are:
1. A positive stochastic process can be defined on the future mass-velocity bundle.
2. Timelike position trajectories can be maintained.
3. The invariant mass can be promoted to a positive stochastic variable.
4. The construction has a natural forward/backward stochastic structure.
5. Its spacetime marginal satisfies a probability conservation equation.
What remains unproved is whether the resulting stochastic dynamics generates Fanchi's quadratic wave operator.

I would now perform a specific calculation rather than develop more general formalism.

Take the free particle, with no electromagnetic field, and start with the simplest model:
\[
dX^\mu=\frac{\mu c}{M}U^\mu ds,
\]

\[
dY=a\,ds+\sigma\,dB_s,
\]

\[
dU=\text{Brownian diffusion on }H^3.
\]

Then:
- Derive its exact forward and backward generators.
- Calculate the stochastic mean acceleration.
- Formulate the stochastic variational principle.
- Investigate whether the internal variables can be eliminated.
- Determine whether the resulting spacetime equation is Klein–Gordon, Fanchi's parametrized equation, or something different.
The free-particle case is particularly valuable because there is no ambiguity associated with electromagnetic gauge coupling or spin.
A useful criterion for success

The theory should recover
\[
i\hbar\partial_s\psi
=
-\frac{\hbar^2}{2M}\Box\psi
\]

in the appropriate limit, while retaining a well-defined positive probability distribution and causal particle trajectories.

If it does not, we should establish precisely which of these requirements is incompatible with the chosen stochastic process.

One final observation: I would not impose the fixed mass-shell condition at the outset. Doing so would remove one of the most interesting features of SHP and Fanchi's theories.

Instead, I would investigate whether a distribution of positive invariant masses can arise dynamically, with a sharply defined mass emerging as a stationary or limiting state.

That would connect stochastic mechanics, off-shell relativistic dynamics and the physical origin of particle mass in a particularly interesting way.
Let us carry out the proposed calculation explicitly. The first question is whether a relativistic diffusion on the space of positive masses and future-directed four-velocities can be made mathematically consistent, and what sort of spacetime dynamics it induces.

The second, much harder question is whether that dynamics can reproduce Fanchi's parametrized wave equation. We should keep these questions separate: constructing a relativistic stochastic process is not yet deriving quantum mechanics from it.
1. The stochastic process
Use the metric
\[
\eta_{\mu\nu}=\operatorname{diag}(+,-,-,-).
\]

Write the four-momentum as
\[
P^\mu=\mu c\,U^\mu,
\qquad
U_\mu U^\mu=1,\quad U^0>0,\quad \mu>0.
\]

Here \(\mu\) is the instantaneous invariant mass and \(U^\mu\) is the four-velocity normalized independently of that mass. The velocity variable takes values on the future unit hyperboloid
\[
H^3=\{U\in\mathbb R^{1,3}:U_\mu U^\mu=1,\ U^0>0\}.
\]

Its induced positive-definite Riemannian metric is the negative of the restriction of the Minkowski metric to its tangent spaces.

For the mass variable introduce
\[
Y=\ln(\mu/\mu_*),\qquad \mu=\mu_*e^Y.
\]

Consider the following process, parametrized by the invariant parameter \(s\):
\[
\begin{aligned}
dX^\mu_s&=V^\mu(Y_s,U_s)\,ds,\\
dY_s&=a(Y_s)\,ds+\sigma\,dB_s,\\
dU_s&=b(U_s)\,ds+\sqrt{\nu}\,dW_s^{H^3}.
\end{aligned}
\tag{1}
\]

The last equation is understood geometrically: \(W^{H^3}\) denotes Brownian motion on the hyperboloid, with the drift \(b\) tangent to it. The coefficients \(a,\sigma,\nu\) may be taken constant initially.

The spacetime velocity is
\[
V^\mu(Y,U)=\frac{\mu c}{M}U^\mu
=\frac{\mu_*c}{M}e^YU^\mu.
\tag{2}
\]

Consequently,
\[
V_\mu V^\mu=\frac{\mu^2c^2}{M^2}>0,\qquad V^0>0.
\tag{3}
\]

This is an important feature: every sample path \(X_s\) is future-directed and timelike, even though its velocity is itself stochastic. The position has finite variation, rather than a Wiener increment. We avoid the familiar difficulty of trying to construct a positive probability distribution for a four-dimensional Wiener process with covariance proportional to the indefinite Minkowski metric.

The price is that the stochasticity is in the mass and four-velocity variables, not directly in spacetime position.
2. Forward generator and Fokker–Planck equation
The backward generator of (1), acting on functions \(g(x,Y,U)\), is
\[
L_+g
=
V^\mu\partial_\mu g
+a\partial_Yg
+\frac{\sigma^2}{2}\partial_Y^2g
+b^A\nabla_Ag
+\frac{\nu}{2}\Delta_{H^3}g.
\tag{4}
\]

Here \(\nabla_A\) and \(\Delta_{H^3}\) refer to the Riemannian geometry of the unit hyperboloid.

Let \(f(x,Y,U,s)\) be the probability density with respect to
\[
d^4x\,dY\,d\operatorname{vol}_{H^3}.
\]

The forward equation is the adjoint equation:
\[
\begin{aligned}
\partial_s f={}&-\partial_\mu(V^\mu f)
-\partial_Y(af)
+\frac{\sigma^2}{2}\partial_Y^2f\\
&-\operatorname{div}_{H^3}(bf)
+\frac{\nu}{2}\Delta_{H^3}f.
\end{aligned}
\tag{5}
\]

This is a relativistic kinetic Fokker–Planck equation. Its structure is close to the relativistic diffusion equations studied by Haba and Herrmann, although the additional stochastic mass coordinate \(Y\) is an extension of the usual fixed-mass construction. The hyperboloid is precisely what makes the velocity diffusion compatible with a positive probability measure.
3. Spacetime probability and current
Integrating over the internal variables gives
\[
\rho(x,s)=
\int_{\mathbb R\times H^3}f(x,Y,U,s)\,
dY\,d\operatorname{vol}_{H^3}.
\tag{6}
\]

Likewise define
\[
J^\mu(x,s)=
\int_{\mathbb R\times H^3}V^\mu(Y,U)f(x,Y,U,s)\,
dY\,d\operatorname{vol}_{H^3}.
\tag{7}
\]

Assuming sufficient decay at the boundaries of the internal space, integration of (5) yields
\[
\boxed{\partial_s\rho+\partial_\mu J^\mu=0.}
\tag{8}
\]

Thus we have obtained a positive spacetime probability density and a conserved spacetime probability current.

But notice a significant difference from Fanchi's theory. Fanchi writes
\[
J^\mu=\frac{\rho}{M}
\left(\partial^\mu S-\frac qc A^\mu\right),
\tag{9}
\]

whereas our current is an average over a distribution of masses and four-velocities. In general,
\[
J^\mu\ne \rho\,V^\mu_{\rm single}.
\]

There is no reason yet for the current to be the gradient of a single phase.

This is the first major obstacle to recovering Fanchi's equation.
4. Forward and backward drifts
Now introduce the specifically Nelsonian part of the construction.

For a diffusion with density \(f\), the backward internal drifts are
\[
a_-=a-\sigma^2\partial_Y\ln f,
\tag{10}
\]

and
\[
b_-^A=b^A-\nu\nabla^A\ln f.
\tag{11}
\]

The corresponding backward generator is
\[
L_-=
V^\mu\partial_\mu
+a_-\partial_Y+\frac{\sigma^2}{2}\partial_Y^2
+b_-^A\nabla_A+\frac{\nu}{2}\Delta_{H^3}.
\tag{12}
\]

The forward and backward drifts of \(X^\mu\) coincide because \(X^\mu\) has no direct diffusion:
\[
D_+X^\mu=D_-X^\mu=V^\mu.
\tag{13}
\]

The internal osmotic velocities are
\[
u_Y=\frac{\sigma^2}{2}\partial_Y\ln f,
\qquad
u^A=\frac{\nu}{2}\nabla^A\ln f,
\tag{14}
\]

while the corresponding current drifts are
\[
v_Y=a-\frac{\sigma^2}{2}\partial_Y\ln f,
\qquad
v^A=b^A-\frac{\nu}{2}\nabla^A\ln f.
\tag{15}
\]

This is a genuine forward/backward diffusion structure, but its osmotic velocity lives in the internal variables \((Y,U)\), not in spacetime.
5. Nelson's mean acceleration
Define the symmetric stochastic mean acceleration by
\[
\mathcal A^\mu
=\frac12\left(D_+D_-+D_-D_+\right)X^\mu.
\tag{16}
\]

Since \(D_+X^\mu=D_-X^\mu=V^\mu\), this becomes
\[
\mathcal A^\mu=\frac12(L_++L_-)V^\mu.
\tag{17}
\]

We can calculate this explicitly.

Set
\[
C=\frac{\mu_*c}{M},\qquad V^\mu=Ce^YU^\mu.
\]

The derivatives with respect to \(Y\) are particularly simple:
\[
\partial_YV^\mu=V^\mu,\qquad
\partial_Y^2V^\mu=V^\mu.
\tag{18}
\]

For the hyperboloid, with the convention
\(\Delta=\operatorname{div}\operatorname{grad}\), its embedding coordinates satisfy
\[
\Delta_{H^3}U^\mu=3U^\mu.
\tag{19}
\]

Also,
\[
\nabla_A V^\mu=Ce^Y\nabla_AU^\mu.
\]

Substitution into (17) gives
\[
\boxed{
\begin{aligned}
\mathcal A^\mu={}&
\left[
a+\frac{\sigma^2}{2}+\frac{3\nu}{2}
-\frac{\sigma^2}{2}\partial_Y\ln f
\right]V^\mu\\
&+\left[
b^A-\frac{\nu}{2}\nabla^A\ln f
\right]\nabla_A V^\mu.
\end{aligned}}
\tag{20}
\]

This is a useful result. The mean acceleration has two parts:
1. A component parallel to \(V^\mu\), arising from the mass drift, mass diffusion, and hyperbolic diffusion.
2. A component tangent to the velocity hyperboloid, arising from the drift and osmotic velocity on \(H^3\).
The density-gradient terms are the relativistic kinetic analogue of Nelson's osmotic contribution.

For the simplest free model, take \(a=b=0\). Then
\[
\mathcal A^\mu
=
\left[
\frac{\sigma^2+3\nu}{2}
-\frac{\sigma^2}{2}\partial_Y\ln f
\right]V^\mu
-\frac{\nu}{2}\nabla^A\ln f\,\nabla_A V^\mu.
\tag{21}
\]

Even in the absence of external forces, the mean acceleration need not vanish. This is not a contradiction: the process is diffusive in its internal variables.
A point about the mass shell
For a fixed-mass relativistic particle, the four-acceleration is orthogonal to the four-velocity:
\[
U_\mu\frac{dU^\mu}{ds}=0.
\]

But in our model the mass fluctuates, and
\[
V_\mu V^\mu=\frac{\mu^2c^2}{M^2}.
\]

Consequently,
\[
V_\mu\mathcal A^\mu
=\frac12\left(D_++D_-\right)(V_\mu V^\mu)/2
\]

with the precise stochastic derivative understood through the generator. There is no general requirement that this contraction vanish.

This matters when we try to introduce the electromagnetic Lorentz force. The ordinary Lorentz force is orthogonal to the four-velocity:
\[
F_{\mu\nu}V^\mu V^\nu=0.
\]

It therefore cannot, by itself, account for stochastic changes in invariant mass. A mass-changing model needs an additional longitudinal interaction or a separate stochastic mass law.
6. The spacetime moment hierarchy
There is another way to see why this construction does not immediately produce Fanchi's equation.

Define the second spacetime moment
\[
\Pi^{\mu\nu}(x,s)
=
\int V^\mu V^\nu f\,
dY\,d\operatorname{vol}_{H^3}.
\tag{22}
\]

Multiply the Fokker–Planck equation (5) by \(V^\mu\) and integrate over the internal variables. Integration by parts gives
\[
\boxed{
\partial_sJ^\mu+\partial_\nu\Pi^{\mu\nu}
=
\int (L_{\rm int}V^\mu)f\,
dY\,d\operatorname{vol}_{H^3},
}
\tag{23}
\]

where
\[
L_{\rm int}
=a\partial_Y+\frac{\sigma^2}{2}\partial_Y^2
+b^A\nabla_A+\frac{\nu}{2}\Delta_{H^3}.
\tag{24}
\]

Equation (23) is the analogue of a momentum-balance equation in kinetic theory.

But it contains \(\Pi^{\mu\nu}\), and the equation for \(\Pi^{\mu\nu}\) will involve third moments. We therefore obtain a hierarchy:
\[
\rho,\quad J^\mu,\quad \Pi^{\mu\nu},\quad
\text{third moments},\ldots
\]

There is no automatic closure at the level of \(\rho\) and \(J^\mu\).

This is a fundamental difference between a kinetic stochastic theory and a wavefunction theory. In Fanchi's equation the current is determined directly by \(\psi\), whereas here the current is a marginal of a distribution on a larger space.
7. Can the Fanchi equation emerge?
The target equation is
\[
i\hbar\partial_s\psi
=
\frac{1}{2M}
\left(-i\hbar\partial_\mu-\frac qc A_\mu\right)
\left(-i\hbar\partial^\mu-\frac qc A^\mu\right)\psi.
\tag{25}
\]

Writing
\[
\psi=\sqrt{\rho}\,e^{iS/\hbar}
\]

gives the two real equations
\[
\partial_s\rho+
\partial_\mu\left[
\frac{\rho}{M}
\left(\partial^\mu S-\frac qcA^\mu\right)
\right]=0,
\tag{26}
\]

and
\[
\partial_sS+
\frac{1}{2M}
\left(\partial_\mu S-\frac qcA_\mu\right)
\left(\partial^\mu S-\frac qcA^\mu\right)
-\frac{\hbar^2}{2M}
\frac{\Box\sqrt{\rho}}{\sqrt{\rho}}=0.
\tag{27}
\]

Our stochastic model has produced an equation of the form (26), but only in the weaker sense of a marginal continuity equation. It has not produced the particular current in (26), and it has not produced the quantum potential
\[
Q=-\frac{\hbar^2}{2M}\frac{\Box\sqrt{\rho}}{\sqrt{\rho}}.
\tag{28}
\]

In particular, the diffusion in \(Y\) and \(U\) naturally produces derivatives with respect to internal variables. It does not automatically produce the spacetime d'Alembertian \(\Box\).

To get (27), we would need an additional result: after eliminating the internal variables, the stochastic variational principle must generate precisely the spacetime quantum potential (28). It would be incorrect simply to insert that term into the equations.
8. What does the existing literature tell us?
There are two relevant, but distinct, strands of work.
- Marra and Serva developed variational principles for relativistic stochastic mechanics and obtained Klein–Gordon dynamics under appropriate assumptions. Their paper is particularly relevant because it treats the stochastic variational problem, rather than merely writing down a relativistic diffusion. See Annales de l'Institut Henri Poincaré A 53 (1990), 97–108.
- Haba and Herrmann studied relativistic diffusions on the velocity hyperboloid, including external forces and kinetic equations. These works provide much of the geometric and probabilistic machinery needed for the internal diffusion proposed here. See Haba, Physical Review E 79 (2009), 021128, and Herrmann, Physical Review E 80 (2009), 051110.
These results establish that relativistic stochastic dynamics is not an empty programme. But they do not establish that the particular fluctuating-mass model (1) yields Fanchi's parametrized wave equation.

9. My assessment: where the real problem lies
I would distinguish three levels of achievement:
Level	What is required	Status of our construction
Relativistic stochastic kinematics	Positive probability, causal spacetime paths	Achieved
Relativistic stochastic dynamics	A physically justified drift and variational principle	Not yet achieved
Quantum dynamics	Derivation of the Fanchi/Klein–Gordon equation and its quantum potential	Not achieved

The most important issue is not the geometry of \(H^3\). That part is manageable. The real issue is the relation between the internal stochastic degrees of freedom and the spacetime quantum potential.

There is also a conceptual choice to make. We have introduced stochastic mass as an independent degree of freedom. But Fanchi's off-shell dynamics does not necessarily require a stochastic mass: in the free theory, the invariant mass squared can remain constant along a classical trajectory while different quantum mass sectors coexist in a superposition.

Thus, stochastic mass fluctuations may be an unnecessary complication.

I suggest the next calculation should compare two models:
1. A fixed-mass model, with diffusion only on \(H^3\).
2. A fluctuating-mass model, with diffusion on \(\mathbb R\times H^3\).
For each, we should derive the stochastic variational equations and ask whether the marginal spacetime dynamics closes. The fixed-mass model is the cleaner starting point, and it will tell us whether the extra mass coordinate contributes anything essential.

My preference is to start with the fixed-mass model. It is closer to the relativistic diffusion literature, avoids adding an unmotivated stochastic degree of freedom, and provides a sharper test of whether Nelson's stochastic mechanism can produce Fanchi's parametrized quantum dynamics.
Let us now take the fixed-mass model and see what it actually gives. This is a useful test because it removes the extra stochastic mass variable and leaves us with the geometry of relativistic velocity space.

There is, however, an important distinction to maintain. A fixed-mass relativistic diffusion is naturally parametrized by proper time. Fanchi's theory uses an independent invariant parameter \(s\), and its quantum states need not be restricted to one mass sector. We can compare the two, but we should not identify them at the outset.
1. Fixed-mass relativistic diffusion
Let
\[
P^\mu=mc\,U^\mu,\qquad U_\mu U^\mu=1,\quad U^0>0.
\]

The velocity variable belongs to \(H^3\), the future unit hyperboloid.

Use proper time \(\tau\) as the evolution parameter. The spacetime trajectory satisfies
\[
dX^\mu_\tau=cU^\mu_\tau\,d\tau.
\tag{1}
\]

Thus every sample path is future-directed and timelike.

For the simplest free diffusion, take Brownian motion on \(H^3\):
\[
dU_\tau=\sqrt{\nu}\,dW^{H^3}_\tau.
\tag{2}
\]

Here \(\nu\) has dimensions of inverse time if \(\tau\) is measured in time units. More generally, a physically realistic process may include a tangent drift \(b\), for example a drag term.

The backward generator is
\[
L_+=cU^\mu\partial_\mu+\frac{\nu}{2}\Delta_{H^3}.
\tag{3}
\]

The corresponding density \(f(x,U,\tau)\), with respect to
\(d^4x\,d\operatorname{vol}_{H^3}\), satisfies
\[
\boxed{
\partial_\tau f+
cU^\mu\partial_\mu f
=\frac{\nu}{2}\Delta_{H^3}f.
}
\tag{4}
\]

This is a relativistic kinetic diffusion equation. It is not yet a wave equation.
2. Forward and backward velocities
For a density \(f\), the backward drift on the hyperboloid is
\[
b_-^A=-\nu\nabla^A\ln f,
\tag{5}
\]

when the forward drift is zero. Hence the osmotic velocity is
\[
u^A=\frac{\nu}{2}\nabla^A\ln f,
\tag{6}
\]

and the current velocity is
\[
v^A=-\frac{\nu}{2}\nabla^A\ln f.
\tag{7}
\]

The position variable has finite variation:
\[
D_+X^\mu=D_-X^\mu=cU^\mu.
\tag{8}
\]

Consequently, the spacetime osmotic velocity is zero. All the Nelsonian osmotic structure is in the velocity variables.

This is already a substantial departure from ordinary Nelson mechanics, where the configuration variable itself is diffusive.
3. Mean acceleration on the hyperboloid
The calculation can be made explicit. For the forward generator,
\[
L_+=cU^\mu\partial_\mu+\frac{\nu}{2}\Delta_{H^3},
\]

we use the identity
\[
\Delta_{H^3}U^\mu=3U^\mu.
\tag{9}
\]

The backward generator is
\[
L_-=cU^\mu\partial_\mu
-\nu\nabla^A\ln f\,\nabla_A
+\frac{\nu}{2}\Delta_{H^3}.
\tag{10}
\]

Applying the two generators to \(cU^\mu\), the symmetric mean acceleration becomes
\[
\boxed{
\mathcal A^\mu
=
\frac{3\nu c}{2}U^\mu
-\frac{\nu c}{2}
\nabla^A\ln f\,\nabla_AU^\mu .
}
\tag{11}
\]

The second term is tangent to \(H^3\), and hence orthogonal to \(U^\mu\):
\[
U_\mu\nabla_AU^\mu=0.
\tag{12}
\]

The first term is parallel to \(U^\mu\). This is not an ordinary classical four-acceleration. It is a feature of the symmetric stochastic mean acceleration; the stochastic velocity has fluctuations even though its value remains on the unit hyperboloid.

In particular, the limit \(\nu\to0\) gives
\[
\mathcal A^\mu\longrightarrow0,
\]

as expected for free inertial motion.

One should not interpret (11) as a new force law without further work. It is a kinematical consequence of the chosen diffusion.
4. The moment equations: where closure fails
Define the spacetime marginal density and its first two moments:
\[
\begin{aligned}
\rho(x,\tau)&=\int_{H^3}f\,dU,\\
J^\mu(x,\tau)&=c\int_{H^3}U^\mu f\,dU,\\
\Pi^{\mu\nu}(x,\tau)&=
c^2\int_{H^3}U^\mu U^\nu f\,dU.
\end{aligned}
\tag{13}
\]

Here \(dU=d\operatorname{vol}_{H^3}\).

Integrating (4) gives
\[
\boxed{\partial_\tau\rho+\partial_\mu J^\mu=0.}
\tag{14}
\]

Multiplying by \(cU^\mu\), integrating, and using the self-adjointness of the hyperbolic Laplacian gives
\[
\boxed{
\partial_\tau J^\mu+
\partial_\nu\Pi^{\mu\nu}
=\frac{3\nu}{2}J^\mu.
}
\tag{15}
\]

This equation is particularly instructive. The current is not determined by the density alone. Its evolution involves the second moment \(\Pi^{\mu\nu}\). The equation for that moment involves third moments, and so on.

Thus the kinetic description has an infinite moment hierarchy.

We could impose a closure approximation, for example
\[
\Pi^{\mu\nu}\approx
\frac{J^\mu J^\nu}{\rho},
\tag{16}
\]

but that is a monokinetic approximation. It suppresses the very velocity dispersion that makes the process stochastic. It is therefore not a satisfactory route to a quantum theory.

This is the central mathematical obstruction in the fixed-mass model.
5. Relation to Fanchi's parameter \(s\)
There is nevertheless a direct classical connection.

Fanchi's free parametrized Hamiltonian is
\[
K=\frac{1}{2M}p_\mu p^\mu.
\tag{17}
\]

Hamilton's equation gives
\[
\frac{dX^\mu}{ds}=\frac{p^\mu}{M}.
\tag{18}
\]

On the fixed mass shell,
\[
p^\mu=mcU^\mu,\qquad p_\mu p^\mu=m^2c^2.
\tag{19}
\]

Choosing \(M=m\), we obtain
\[
\frac{dX^\mu}{ds}=cU^\mu.
\tag{20}
\]

Since \(U^\mu= dX^\mu/(c\,d\tau)\), this means
\[
\boxed{ds=d\tau.}
\tag{21}
\]

So the classical fixed-mass sector of Fanchi's parametrized dynamics is precisely proper-time evolution.

This is a useful identification, but it also exposes a limitation: a fixed-mass hyperboloid diffusion explores only the mass shell. It cannot by itself describe Fanchi's full off-shell state space, in which different values of \(p_\mu p^\mu\) may occur.

The quantum version makes the distinction especially clear. Fanchi's free equation is
\[
i\hbar\partial_s\psi
=-\frac{\hbar^2}{2M}\Box\psi.
\tag{22}
\]

A stationary mode,
\[
\psi(x,s)=\phi(x)e^{-i\kappa s/\hbar},
\]

satisfies
\[
-\hbar^2\Box\phi=2M\kappa\,\phi.
\tag{23}
\]

The Klein–Gordon sector of rest mass \(m\) corresponds to
\[
2M\kappa=m^2c^2.
\tag{24}
\]

Thus a single mass-shell sector is a stationary \(s\)-mode of the parametrized theory. A superposition of such sectors is more general than a single fixed-mass diffusion.
6. Can the diffusion coefficient be related to \(\hbar\)?
In ordinary Nelson mechanics, the diffusion coefficient is
\[
\nu_{\rm N}=\frac{\hbar}{2m},
\tag{25}
\]

which has dimensions of length squared divided by time. It is the diffusion coefficient of the configuration variable.

Here, however, the Brownian variable is the dimensionless four-velocity \(U^\mu\). Consequently, \(\nu\) has dimensions of inverse time. We cannot simply identify it with \(\hbar/(2m)\).

There is a natural relativistic scale:
\[
\omega_C=\frac{mc^2}{\hbar},
\tag{26}
\]

the Compton frequency. We could therefore write
\[
\nu=\alpha\frac{mc^2}{\hbar},
\tag{27}
\]

where \(\alpha\) is dimensionless.

But this is only dimensional reasoning. Nothing in the diffusion model determines \(\alpha\), let alone proves that it has the value needed to recover the quantum theory.

This is an important difference from ordinary Nelson mechanics, where the relation between the diffusion coefficient and \(\hbar\) is central to the derivation of Schrödinger dynamics.
7. Comparison with relativistic stochastic mechanics in the literature
The fixed-mass hyperboloid construction is not the only possible relativistic stochastic mechanics.

Marra and Serva, in their 1990 paper Variational principles for a relativistic stochastic mechanics, considered relativistic covariant diffusion processes and developed stochastic variational principles. Under appropriate conditions, they obtained the Klein–Gordon equation. Their work is particularly important because it addresses the dynamical variational problem, rather than merely constructing a relativistic probability process.

Haba's and Herrmann's work on relativistic diffusion provides a complementary approach: the diffusion takes place on the velocity hyperboloid, with corresponding relativistic kinetic equations.

These approaches should not be conflated:
- The hyperboloid diffusion gives a well-defined stochastic process with causal trajectories.
- The stochastic variational approach aims to derive relativistic wave dynamics.
- Fanchi's parametrized theory gives a unitary evolution in an invariant parameter and allows off-shell states.
The interesting possibility is that these three ingredients could be combined, but the combination is not automatic. In particular, the kinetic equation (4) is not the parametrized wave equation (22).

For references:
- R. Marra and M. Serva, “Variational principles for a relativistic stochastic mechanics,” Annales de l'Institut Henri Poincaré A 53 (1990), 97–108.
- Z. Haba, “Relativistic diffusion,” Physical Review E 79 (2009), 021128.
- J. Herrmann, “Diffusion in the special theory of relativity,” Physical Review E 80 (2009), 051110.
8. Where I think we should go next
We have now established two things:
1. A fixed-mass relativistic diffusion on \(H^3\) is geometrically consistent and gives causal spacetime trajectories.
2. Its spacetime marginal does not close, and the Nelsonian mean acceleration does not by itself yield Fanchi's quantum potential.
I would therefore not try to force the kinetic equation into the form of a wave equation. Instead, I would investigate the stochastic variational principle.

There are two possible routes:

Route A — relativistic Nelson variational mechanics.

Start with a fixed-mass process on the hyperboloid, introduce forward and backward stochastic derivatives, and construct an action whose variation gives the stochastic equations of motion. Then determine whether the resulting equations can be reduced to Klein–Gordon dynamics.

Route B — parametrized stochastic mechanics.

Start from Fanchi's classical Hamiltonian
\[
K=\frac{1}{2M}
\left(p_\mu-\frac qc A_\mu\right)
\left(p^\mu-\frac qc A^\mu\right),
\tag{28}
\]

and construct a stochastic variational principle directly in the extended phase space \((X^\mu,p_\mu)\), with \(s\) as evolution parameter. The question is whether its quantum correction produces
\[
Q=-\frac{\hbar^2}{2M}
\frac{\Box\sqrt{\rho}}{\sqrt{\rho}}.
\tag{29}
\]

I favor Route B for our particular purpose. It preserves the essential SHP/Fanchi structure from the beginning, rather than trying to reconstruct it from a proper-time diffusion.

There is one important qualification: the stochastic process must be chosen so that its position component remains causal, or else we must explain physically why individual stochastic paths need not be causal even when observable propagation is relativistically consistent.

The decisive test will be whether a stochastic action produces (29), rather than whether we can write down a diffusion equation that resembles a relativistic equation.

